% ================================================
% 文件: 04_仪器仪表.tex
% 章节: 第17章 仪器仪表
% 所属部分: 系统与设备
% 版本: v1.0
% ================================================
% 本章目录结构（树状）:
% \chapter{仪器仪表}
% ├── \section{仪器仪表概述}
% ├── \section{测量仪器原理}
% ├── \section{常用电子测量仪器}
% ├── \section{仪器仪表校准}
% └── \section{仪器仪表应用}
% ================================================

\chapter{仪器仪表}
\label{chap:instruments}

\section{仪器仪表概述}
\label{sec:instruments_intro}

仪器仪表是用于测量、监测和控制各种物理量的设备。在电子工程领域，仪器仪表是设计、调试和维护电子设备不可或缺的工具。

从工程方法论的角度看，仪器仪表的作用在于把人的感官无法直接感知的物理量（电压、电流、频率、相位、功率、噪声、失真等）转换为可读数、可记录、可比较的量值。通信设备的电气参数大多是高频的、微弱的或瞬态的，既不能用肉眼观察，也无法凭经验判断，因此设备的设计定型、生产验收、安装开通与日常维护，每一个环节都必须以测量数据为依据。可以说，没有可信的测量，就没有可信的技术指标；测量能力的高低，直接决定了一个实验室或维修部门能够承担的工作深度。

一台完整的测量仪器通常由若干功能环节串接而成：\textbf{传感与变换环节}把被测量取出并转换为电信号（如探头、取样电阻、定向耦合器、天线）；\textbf{信号调理环节}完成放大、衰减、滤波、检波、混频等处理，使信号进入后级电路的最佳工作范围；\textbf{测量与运算环节}完成比较、计数、模数转换与数值计算；\textbf{显示与输出环节}以指针、数码、图形或数字接口的形式给出结果。理解这一通用结构，对判断仪器的能力边界与故障部位都很有帮助——例如某台数字电压表读数偏低，故障既可能出在输入衰减网络，也可能出在参考电压源，而不必先怀疑显示电路。

仪器仪表的分类：
\begin{itemize}
    \item \textbf{按功能}：测量仪器、分析仪器、控制仪器
    \item \textbf{按原理}：模拟仪器、数字仪器、智能仪器
    \item \textbf{按用途}：实验室仪器、工业仪器、便携式仪器
\end{itemize}

上述三种分类分别对应三个不同的观察角度。按功能划分时，测量仪器只回答"是多少"，如电压表、频率计、功率计；分析仪器进一步回答"由什么组成"，如频谱分析仪、失真度分析仪、逻辑分析仪；控制仪器则把测量结果反馈到被控对象，构成闭环，如程控电源、恒温控制器。按原理划分时，模拟仪器以连续量的偏转或比较为基础，读数直观、无采样延迟，但精度受机械与电路线性度限制；数字仪器以模数转换和数字计数为基础，读数客观、分辨力高、便于存储与传输；智能仪器在数字仪器基础上引入微处理器，具备自动量程、自校准、数据处理与程控接口的能力。按用途划分时，实验室仪器追求准确度与功能完备，工业仪器强调可靠性与抗环境干扰，便携式仪器则以电池供电、体积小、耐振动为设计目标，其准确度指标通常低于同类台式仪器，选用时应予注意。

讨论仪器性能时必须使用统一的术语，否则容易把不同含义的指标混为一谈。下表列出了几个最常用的基本术语。其中特别需要区分"准确度"与"精密度"：前者描述测量结果接近真值的程度，后者只描述多次测量彼此接近的程度。一台零点漂移的仪器可能重复性极好（精密度高）而读数普遍偏离真值（准确度低）；反之，一台读数分散但无系统偏差的仪器，取多次平均后仍可给出较准确的结果。

\begin{table}[htbp]
\centering
\caption{测量仪器常用基本术语}
\begin{tabular}{|p{2.6cm}|p{10.4cm}|}
\hline
\textbf{术语} & \textbf{含义说明} \\
\hline
量程 & 仪器在规定准确度内所能测量的被测量范围，由下限和上限确定 \\
\hline
分辨力 & 仪器能够区分的被测量最小变化量；数字仪器通常表现为末位一个字 \\
\hline
灵敏度 & 输出变化量与引起该变化的输入变化量之比，反映仪器对微小量的响应能力 \\
\hline
准确度 & 测量结果与被测量真值的接近程度，常以等级或相对误差限表示 \\
\hline
精密度 & 相同条件下多次重复测量结果彼此接近的程度，用标准偏差描述 \\
\hline
稳定性 & 在规定时间内仪器计量特性保持不变的能力，含零点漂移与量程漂移 \\
\hline
输入阻抗 & 仪器输入端对被测电路呈现的等效阻抗，决定接入对被测电路的影响 \\
\hline
\end{tabular}
\end{table}

灵敏度可用输入输出变化量之比定量描述，对于线性刻度的仪器有
\[
S = \frac{\Delta y}{\Delta x},
\]
式中 $\Delta x$ 为被测量的变化量，$\Delta y$ 为相应的输出（指针偏转格数或数码读数）变化量。灵敏度并非越高越好：灵敏度提高往往意味着通带内噪声同样被放大，若仪器本身的噪声电平已与被测量相当，再提高灵敏度只会使读数抖动加剧而不会带来有效信息。因此工程上评价微弱信号测量能力时，更常用"最小可测量"或等效输入噪声来表述，而不是单看灵敏度数值。

\section{测量仪器原理}
\label{sec:measurement_principle}

测量仪器的基本原理是将被测物理量转换为可观测的信号。

测量的实质是把被测量与同类的计量单位相比较，求出比值。这一过程可写为
\[
x = N \cdot u,
\]
式中 $x$ 为被测量，$u$ 为所用的计量单位，$N$ 为比值即读数。因此任何测量结果都由数值和单位两部分组成，缺少单位的读数没有意义；同时也说明，测量的准确度不可能高于所依据单位（即标准量）的准确度，这正是量值溯源的必要性所在。按获得结果的途径不同，测量可分为直接测量（用仪器直接读出被测量，如用电压表测电压）、间接测量（测出与被测量有确定函数关系的其他量后计算得出，如由电压和电阻计算功率）与组合测量（联立求解多个方程得到多个未知量）。间接测量的误差由各直接测量量的误差按函数关系合成，往往大于单项误差，设计测量方案时应尽量减少中间环节。

常见的测量原理：
\begin{itemize}
    \item \textbf{电桥原理}：用于电阻、电容、电感测量
    \item \textbf{放大原理}：用于微弱信号测量
    \item \textbf{比较原理}：与标准量比较进行测量
    \item \textbf{数字化原理}：将模拟信号转换为数字信号
\end{itemize}

\textbf{电桥原理}属于零示法的典型应用。以直流单臂电桥测电阻为例，四个桥臂分别为 $R_1$、$R_2$、$R_3$ 和待测电阻 $R_x$，调节可调臂使检流计指零，此时电桥平衡，满足
\[
R_1 R_x = R_2 R_3, \qquad R_x = \frac{R_2 R_3}{R_1}.
\]
可见待测值仅由三个已知电阻决定，与电源电压和检流计的灵敏度标定无关，只要求检流计足够灵敏以判断"零"。这使电桥法能获得很高的准确度。交流电桥用于测量电容和电感时，平衡条件须同时满足阻抗模与相角两个方程，因此桥臂中一般需要两个可调元件，并要求电源波形纯净、屏蔽良好，否则分布电容与外界干扰会使平衡点漂移。

\textbf{放大原理}用于被测量太小、不足以驱动指示机构的场合。放大环节在提升信号幅度的同时也放大了自身的噪声与漂移，因此微弱信号测量的关键不在增益而在信噪比。工程上常用的对策有：采用低噪声输入级并使源阻抗与之匹配；用调制（斩波）方式把直流信号搬移到较高频率以避开 $1/f$ 噪声；用带通滤波或相关检测压缩等效噪声带宽。由于噪声功率与带宽成正比，等效噪声电压与带宽的平方根成正比，即
\[
U_n = \sqrt{4kTRB},
\]
式中 $k$ 为玻尔兹曼常数，$T$ 为绝对温度，$R$ 为等效噪声源电阻，$B$ 为等效噪声带宽。这说明只要允许测量时间足够长（带宽足够窄），微弱信号的可测下限就能显著降低。

\textbf{比较原理}又可细分为零示法、微差法与替代法。零示法如上述电桥，用指零仪判断平衡；微差法把被测量与标准量之差送入仪表，只要求仪表能准确测量这一小量，因而总误差被大幅压缩；替代法先用被测量使仪器得到某一读数，再用可调标准量替代被测量使读数复原，此时标准量即等于被测量，仪器本身的刻度误差被完全消除，只要求仪器读数具有良好的重复性。射频功率与衰减量的精密测量大量使用替代法。

\textbf{数字化原理}是现代仪器的基础。模数转换包含采样、量化、编码三步。采样必须满足采样定理，即采样频率不低于被测信号最高频率的两倍
\[
f_s \ge 2 f_{\max},
\]
否则将产生频谱折叠（混叠），在屏幕上表现为频率虚假降低的波形，因此实际仪器都在采样前设置抗混叠低通滤波器。量化把连续幅度归并到有限个电平，$n$ 位转换器的量化台阶为
\[
q = \frac{U_{FS}}{2^{n}},
\]
最大量化误差为 $\pm q/2$，理想情况下信噪比约为 $6.02n + 1.76$~dB。可见提高分辨位数可以降低量化误差，但转换器的实际有效位数往往低于标称位数，还受参考源稳定度、线性误差和输入噪声的限制，选用时应关注有效位数与积分非线性等指标。

测量误差来源：系统误差、随机误差、粗大误差。

误差是测量结果与真值之差。绝对误差定义为 $\Delta = x - A_0$，其中 $x$ 为测量值，$A_0$ 为真值（实际工作中以约定真值或更高等级标准的示值代替）。绝对误差不便于比较不同量级的测量质量，故常用相对误差
\[
\gamma = \frac{\Delta}{A_0} \times 100\%.
\]
对于具有指示刻度的仪表，还使用引用误差（满度相对误差）
\[
\gamma_m = \frac{\Delta_{\max}}{x_m} \times 100\%,
\]
式中 $x_m$ 为满量程值，$\Delta_{\max}$ 为量程内的最大绝对误差。仪表的准确度等级即按引用误差划分，例如 1.0 级表示引用误差不超过 $1.0\%$。由此可导出一条重要的使用规则：由于同一量程内 $\Delta_{\max}$ 大体固定，读数越接近满度，相对误差越小。因此测量时应尽量选择使指针指在满刻度三分之二以上的量程，避免在小于三分之一量程处读数。

三类误差的性质与处理方式差别很大，见下表。系统误差在重复测量中保持恒定或按确定规律变化，取平均值不能减小，只能通过修正值、改进方法或改换标准来消除，它决定了测量的准确度；随机误差由大量微小、互不相关的因素造成，符合统计规律，可通过增加测量次数取平均来减小，它决定了测量的精密度；粗大误差明显超出正常范围，多由读错、记错、操作失误或外界剧烈干扰造成，应按判别准则剔除后再作数据处理。

\begin{table}[htbp]
\centering
\caption{三类测量误差的特征与处理}
\begin{tabular}{|p{2.0cm}|p{4.4cm}|p{3.0cm}|p{3.4cm}|}
\hline
\textbf{类别} & \textbf{主要来源} & \textbf{表现特征} & \textbf{处理方法} \\
\hline
系统误差 & 仪器刻度偏差、零点漂移、方法不完善、环境条件偏离标准值 & 恒定或按规律变化，符号与大小可预期 & 加修正值、改换测量方法、校准与溯源 \\
\hline
随机误差 & 噪声、接触电阻波动、读数视差、微小振动 & 时正时负，多次测量近似正态分布 & 多次测量取算术平均，用标准偏差评定 \\
\hline
粗大误差 & 读数或记录错误、误接线、瞬时强干扰 & 明显偏离其余数据，孤立异常 & 按判别准则剔除，必要时重新测量 \\
\hline
\end{tabular}
\end{table}

对同一量作 $n$ 次等精度测量，最佳估计值取算术平均值 $\bar{x}$，单次测量的实验标准偏差按贝塞尔公式计算
\[
s = \sqrt{\frac{1}{n-1}\sum_{i=1}^{n}(x_i - \bar{x})^2},
\]
而平均值的标准偏差为 $s_{\bar{x}} = s/\sqrt{n}$。后一式说明增加测量次数可以减小随机误差的影响，但收益按 $\sqrt{n}$ 递减，把测量次数从 4 次增加到 16 次只能使不确定度减半，因此实际工作中通常取 5 至 10 次即可，继续增加次数不如设法消除系统误差来得有效。

\section{常用电子测量仪器}
\label{sec:common_instruments}

常用的电子测量仪器：

\begin{itemize}
    \item \textbf{示波器}：观察电信号波形，测量电压、频率、周期
    \item \textbf{信号发生器}：产生各种波形的电信号
    \item \textbf{万用表}：测量电压、电流、电阻等基本电量
    \item \textbf{频谱分析仪}：分析信号的频率成分
    \item \textbf{网络分析仪}：测量网络的S参数、阻抗匹配
\end{itemize}

\textbf{万用表与电压测量。}万用表是使用最广的通用仪表，可测直流与交流电压、电流、电阻，多数还兼有二极管、通断与电容挡。指针式万用表的直流电压挡由表头串联倍压电阻构成，其输入阻抗以"欧姆每伏"表示，低灵敏度表头接入高阻电路时会因分流使读数明显偏低，这是使用指针表最常见的误判来源。数字万用表以模数转换为核心，直流电压挡输入阻抗通常很高，接入影响小，读数以"$\pm$（读数的百分数 $+$ 若干字）"的形式给出准确度。交流测量还须区分响应方式：平均值响应的仪表按正弦波刻度，对正弦波读数正确，其波形因数为
\[
K_f = \frac{U}{\overline{|u|}} \approx 1.11,
\]
一旦被测信号偏离正弦（如脉冲、方波、削顶的射频包络），读数将产生可观误差；真有效值响应的仪表直接按定义
\[
U = \sqrt{\frac{1}{T}\int_{0}^{T} u^2(t)\,\mathrm{d}t}
\]
测量，可用于非正弦波形，但受峰值因数与带宽限制。测量通信设备中的脉冲或数字信号时，应优先选用真有效值仪表并核对其峰值因数指标。

\textbf{频率与时间测量。}电子计数器是频率测量的主要手段，其原理是在标准闸门时间 $T$ 内对被测信号计数，得到 $N$ 个脉冲，则
\[
f_x = \frac{N}{T}.
\]
由于计数起止与被测信号相位无关，计数值存在 $\pm 1$ 个字的固有误差，测频的相对误差为
\[
\frac{\Delta f}{f_x} = \pm\left(\frac{1}{f_x T} + \frac{\Delta f_c}{f_c}\right),
\]
式中第二项为时基（晶振）的相对误差。可见闸门时间越长、被测频率越高，$\pm 1$ 字误差的影响越小；反之测量低频时误差很大，此时应改用测周期方式，即用标准频率对被测信号的一个或多个周期计数，再取倒数。计数器的时基通常采用恒温晶振或温补晶振，其日老化率与温度系数最终决定了长期测量的准确度，因此计数器同样需要定期溯源。

\textbf{示波器。}示波器把电压随时间的变化显示为二维图形，垂直系统决定幅度测量能力，水平系统决定时间测量能力，触发系统决定波形能否稳定重现。垂直通道的带宽定义为幅频响应下降 3~dB 的频率，它与所能观测的最快上升时间近似满足
\[
t_r \approx \frac{0.35}{BW},
\]
对高斯型响应而言此系数约为 0.35。测量脉冲边沿时，示波器显示的上升时间是信号与仪器响应的合成，即
\[
t_{r,\text{显示}} = \sqrt{t_{r,\text{信号}}^2 + t_{r,\text{示波器}}^2},
\]
因此欲使测量误差可忽略，仪器带宽一般应取被测信号最高有效频率的三到五倍。数字示波器还须关注采样率与存储深度：采样率决定单次捕获的时间分辨力，存储深度决定在给定采样率下可记录的时间长度。探头是常被忽视的误差源，$10:1$ 衰减探头在降低容性负载的同时把信号衰减十倍，使用前必须做低频补偿调节，使方波显示平坦；地线过长引入的电感会在快沿波形上产生振铃，测量高速信号时应使用短接地弹簧。

\textbf{频谱分析仪。}频谱分析仪把信号按频率展开，横轴为频率、纵轴为幅度（通常以 dBm 表示），用于观察载波、谐波、杂散、调制边带与噪声。扫频式超外差频谱仪的核心是本振扫频与中频滤波，分辨带宽（RBW）即中频滤波器的带宽，它决定能否分辨两个邻近谱线，也决定仪器的显示平均噪声电平：带宽每减小十倍，噪声电平约下降 10~dB，因此观察微弱杂散时应减小 RBW。代价是扫描时间随 RBW 变窄而迅速增加，近似有
\[
T_{sweep} \propto \frac{\Delta f}{RBW^2},
\]
式中 $\Delta f$ 为扫宽。视频带宽（VBW）用于平滑噪声起伏、使谱线更易读，但不改变噪声的平均电平。使用频谱仪时必须首先估算被测信号功率，必要时前置衰减器或定向耦合器取样，避免输入过载损坏输入混频器——这是频谱仪最常见的损坏原因。

\textbf{网络分析仪。}网络分析仪测量二端口网络的散射参数，即 $S_{11}$、$S_{21}$、$S_{12}$、$S_{22}$。其中 $S_{11}$ 为输入反射系数，与驻波比和回波损耗直接对应；$S_{21}$ 为正向传输系数，反映插入损耗或增益。反射系数与驻波比的关系为
\[
\Gamma = \frac{Z_L - Z_0}{Z_L + Z_0}, \qquad VSWR = \frac{1 + |\Gamma|}{1 - |\Gamma|}.
\]
矢量网络分析仪能同时测量幅度与相位，因而可以通过校准把测试端口的定向性误差、源匹配误差和频响误差数学地扣除，常用的开路—短路—负载—直通校准要求校准件本身量值准确、连接重复性好。校准面即测量参考面，若在校准后加接电缆或转接头，其损耗与相移将被计入被测件，这一点在滤波器和天线馈线测量中尤须注意。

下表归纳了上述仪器的主要用途与选用时应重点核对的指标，可作为配置测量系统时的参考。总的原则是：先明确被测量与所需准确度，再据此确定仪器的量程、带宽与准确度等级，使仪器误差远小于被测参数的允许偏差，通常要求两者相差数倍以上。

\begin{table}[htbp]
\centering
\caption{常用电子测量仪器的用途与关键选用指标}
\begin{tabular}{|p{2.4cm}|p{4.6cm}|p{5.8cm}|}
\hline
\textbf{仪器} & \textbf{主要被测量} & \textbf{关键选用指标} \\
\hline
万用表 & 电压、电流、电阻、通断 & 准确度与位数、输入阻抗、交流响应方式与带宽、安全等级 \\
\hline
电子计数器 & 频率、周期、时间间隔 & 测频上限、时基稳定度与老化率、闸门时间、输入灵敏度 \\
\hline
示波器 & 波形、幅度、时间、边沿、抖动 & 带宽、采样率、存储深度、通道数、触发功能、探头配套 \\
\hline
信号发生器 & 提供激励信号 & 频率范围与准确度、输出电平范围与平坦度、谐波与杂散、调制能力 \\
\hline
频谱分析仪 & 频谱、谐波、杂散、占用带宽 & 频率范围、分辨带宽、显示平均噪声电平、最大安全输入电平 \\
\hline
网络分析仪 & S 参数、驻波比、插损、群延时 & 频率范围、动态范围、端口数、校准件等级 \\
\hline
功率计 & 平均功率、脉冲功率 & 功率量程、频率范围、探头类型、线性度与温度稳定性 \\
\hline
\end{tabular}
\end{table}

\section{仪器仪表校准}
\label{sec:instruments_calibration}

仪器校准是确保测量结果准确可靠的重要环节。

任何仪器的示值都会随元件老化、机械磨损、温湿度变化和搬运振动而缓慢改变，出厂时合格并不意味着长期合格。校准就是在规定条件下，用已知更高准确度的标准去确定被校仪器示值与被测量量值之间的关系，并给出修正值和不确定度。需要注意校准与检定的区别：校准是技术活动，只给出量值关系与不确定度，不判定合格与否；检定是法制计量行为，依据检定规程判定仪器是否符合法定要求。通信设备维修部门日常所做的多为校准，而涉及贸易结算、安全监督或强制管理的计量器具则须依法检定。

校准方法：
\begin{itemize}
    \item \textbf{标准校准}：与国家标准或标准仪器比对
    \item \textbf{自校准}：使用仪器内部标准进行校准
    \item \textbf{定期校准}：按照规定周期进行校准
\end{itemize}

\textbf{标准校准}是量值溯源的基本形式。所谓溯源，是指测量结果能够通过一条具有规定不确定度的、不间断的比较链，与国家计量基准或国际单位制的定义相联系。这条链自上而下依次为国家基准、副基准与工作基准、各级标准器、实验室工作标准，最后到日常使用的仪器；每向下一级传递，不确定度都会增大，因此标准器的准确度一般应比被校仪器高数倍以上，工程上常用测试不确定度比来衡量
\[
TUR = \frac{T}{U},
\]
式中 $T$ 为被校仪器的允许误差限（允差），$U$ 为校准所用方法的扩展不确定度，通常要求 $TUR \ge 4$。若标准与被校仪器准确度相近，校准结果将无法判定被校仪器是否合格。

\textbf{自校准}利用仪器内部的基准源与已知网络，在开机或按键后自动完成零点、增益、通道一致性等项目的修正。它能有效补偿温度漂移和短期漂移，是数字仪器保持标称准确度的必要步骤，例如数字万用表的自零、网络分析仪的用户校准、功率计的调零与参考校准。但自校准依赖机内基准，无法发现基准本身的漂移，因此不能替代外部溯源校准。

\textbf{定期校准}解决的是"多久校一次"的问题。校准周期应根据仪器的稳定性、使用频度、环境条件与测量后果的严重性确定，常用做法是初始按制造商建议或行业惯例设定，之后依据历次校准数据的漂移趋势延长或缩短。周期之间还应安排期间核查，即用稳定的核查标准（如标准电阻、标准信号源）定期比对，以便及早发现异常。校准后应在仪器上粘贴状态标识，一般以合格、限用（部分量程或功能合格）、停用三种状态区分，防止误用。此外，凡遇到仪器过载、跌落、维修或更换关键元件，均应提前校准，不受周期限制。

校准结果的核心是修正值与不确定度。修正值定义为
\[
C = A_0 - x,
\]
即用标准值减去仪器示值，使用时以 $x + C$ 作为修正后的测量结果。各不确定度分量若彼此独立，按方和根合成得到合成标准不确定度
\[
u_c = \sqrt{\sum_{i=1}^{m} u_i^2},
\]
再乘以包含因子得到扩展不确定度 $U = k\,u_c$，工程上通常取 $k = 2$，对应约 $95\%$ 的置信概率。报告测量结果时应写成 $x \pm U$ 并注明包含因子，只给出一个数值而不给出不确定度的结果在计量意义上是不完整的。

校准证书包含：校准日期、校准结果、不确定度、有效期。

除上述四项外，一份规范的校准证书还应载明被校仪器的名称、型号与编号，所用标准器的名称、编号及其溯源信息，校准依据的规程或方法，校准时的环境温度与相对湿度，以及校准人员与核验人员的签署。使用者接到证书后不应只看"合格"二字，而应逐项核对：校准点是否覆盖了自己实际使用的量程与频率；给出的修正值是否需要在日常测量中实际引用；不确定度是否满足本单位测量任务的要求。证书应随仪器建档保存，与历次校准数据一并形成该仪器的计量档案，供分析漂移趋势和调整校准周期使用。

\section{仪器仪表应用}
\label{sec:instruments_applications}

仪器仪表的应用领域：

\begin{itemize}
    \item \textbf{研发测试}：新产品开发、性能测试
    \item \textbf{生产制造}：质量检测、工艺控制
    \item \textbf{维护维修}：故障诊断、设备检修
    \item \textbf{科学研究}：实验测量、数据分析
\end{itemize}

在\textbf{研发测试}阶段，测量的目的是揭示电路的实际行为并与设计预期比较，因此强调仪器功能的完备与测量的灵活性。以通信整机为例，发射支路需测量输出功率、频率准确度、占用带宽、谐波与杂散抑制、调制度或调制误差；接收支路需测量灵敏度、选择性、互调抗扰性与阻塞特性；配套还需测量电源纹波、电流消耗与热性能。这类测量往往需要信号发生器、频谱分析仪、功率计、示波器与假负载等多种仪器组成测量系统，此时应特别注意各仪器之间的阻抗匹配与接地方式，避免测量系统本身成为误差与干扰的来源。

在\textbf{生产制造}阶段，测量的目的是判定产品是否符合规定指标，强调效率、一致性与可判定性。做法上通常把研发阶段的全面测试精简为若干关键项，编成工序检验规程，并借助程控接口实现自动测试；判定时须考虑仪器不确定度与产品允差的关系，必要时采用保护带（把判定限内缩）以降低误判风险。工装夹具的重复性、连接器的插拔寿命与接触电阻，都会直接影响批量测量的一致性，应纳入日常点检。

在\textbf{维护维修}阶段，测量的目的是快速定位故障，强调便携、安全与经验方法的结合。常用的思路是先做整机现象观察与外观检查，再按信号流程分段测量：以电源电压与静态工作点判断偏置是否正常，以示波器或频谱仪逐级观察信号是否存在与是否合格，以替代法或对比法确认可疑单元。射频设备维修时应始终以假负载代替天线，既避免无用辐射，也保护末级功率器件免受失配损伤。测量高压或大功率部位前必须断电泄放，使用符合安全等级的表笔与探头。

在\textbf{科学研究}中，测量往往要求给出量值与不确定度并可复现，因此实验记录必须包含仪器型号、校准状态、环境条件与测量方法，数据处理须说明剔除粗大误差的准则和合成不确定度的方式。这一要求同样适用于工程报告：能被他人重复的测量才具有说服力。

无论用于何种场合，仪器的日常使用与维护都应遵循若干共同要点。使用前核对量程与被测量的量级，先置于较大量程再逐步减小；精密测量前给仪器充分预热，使内部温度稳定；正确选择接地点，避免多点接地造成地环流，射频测量应使用同轴电缆并保持屏蔽连续；测量前后检查探头与电缆是否破损，接头是否清洁并按规定力矩拧紧；含 MOS 器件的仪器与被测板卡应做防静电防护。存放时应置于阴凉干燥通风处，避免长期高温高湿与强磁场附近；长期不用的仪器宜定期通电运行一段时间，以驱除受潮并防止电解电容失效。仪器发生过载、异响、异味或读数异常时，应立即停用并送检，不得带故障继续测量——用不确定其状态的仪器得到的数据，比没有数据更容易导致错误结论。
